Serving a large language model can be divided between two machines: one
processes the prompt and builds the key-value cache of the model, and the
other generates the response from that cache. The cache must then cross a
network, together with a small number of very short messages on which both
machines wait. Data-center networks are designed for this traffic. In this
thesis, we study the same division without that hardware: two GPU servers
connected by commodity Ethernet, with RDMA provided in software and the
cache staged through host memory.

We present libscuffedrdma, a middleware that treats this traffic as mixed.
Small urgent messages and large bulk transfers are given separate endpoints
and separate completion handling behind a single interface, which is
implemented on both libfabric and raw RDMA verbs. On this foundation we build
a credit-based blob channel, a connector for the vLLM serving engine, a
classifier and a rule-tree policy that select a route for each transfer, a
quantizer with a packed wire format, and an NVMe storage tier. The connector
routes by size, and the learned classifier, the rule tree and the
storage tier are evaluated as separate components. We prove the
properties on which the design relies. We also prove that a classifier based on size thresholds cannot recognize decode traffic, and we show that the switching rule of the queue controller contradicts its own derivation.

On the cache of a real model, the quantizer halves the number of bytes
relative to 16-bit floating point, while the eight most-attended positions
agree 91 percent of the time. GPU decompression and storage batching are
measured. On the laboratory pair, the 99th percentile of a 64-byte message
under a stream of 1 MiB writes is \blank{} microseconds with one endpoint per
class and \blank{} microseconds on a shared endpoint, and the transfer of the
cache accounts for \blank{} percent of the time to first token for a prompt
of 2,048 tokens. \waiting{results from E3 and E8} The algorithms are given as pseudocode with their proofs, together with the method by which every number is produced.
